\newcommand{\fontsample}[3][p]{
\begin{figure}[#1]
\centering
\fbox{\begin{minipage}{14.2cm}
\centering
\begin{minipage}{14cm}
\lstinputlisting[language={[LaTeX]{TeX}},style=block,nolol]{code/fonts/header_#2.tex}
\end{minipage}
\includegraphics{figures/samples/fonts/font_#2}
\end{minipage}}
\caption{Font sample "`#3"'}
\label{fig:tpl:font:#2}
\end{figure}}

The typeface has a significant effect on the readability of a text. In
\LaTeX\ three different font families are used for text.

\begin{itemize}
\item {\rmfamily Serif fonts have small horizontal strokes at the ends of the
      strokes, thus emphasising the baseline of a line of text. This way
      lines of text are more perceptible.}
\item {\sffamily Sans-serif fonts do not have these horizontal strokes and are
      therefore more striking. They should therefore be used on title pages,
      posters and in presentations.}
\item {\ttfamily\raggedright The characters of serif typefaces all have the
      same width. Such typefaces are used in mechanical typewriters. They are
      suitable, for example, because of the uniform character size for the
      formatted output of programme code.\par}
\end{itemize}

Furthermore, mathematical fonts complement the characters for formulas.

By default, the font family "`Computer Modern"' is used for all four cases or
its counterpart "Latin Modern"' (see fig.~\ref{fig:tpl:font:lm}) is used for
all application. The latter has a larger character set. However, these fonts
are quite brittle and look very thin when printed on a laser printer.
Publishers therefore often recommend switching to other font families.

\fontsample[h]{lm}{Latin Modern}

\fontsample{newpx}{Palatino}

\fontsample{newtx}{Times Roman}

Loading additional packages redefines these standard font families. Some
packages redefine all families. Others affect only on individual families and
leave the other font families unchanged.

Figure~\ref{fig:tpl:font:newpx} shows which \LaTeX packages are used to change
the standard font to "`Palatino"'. As a sans serif font "Helvetica"' is used.
In addition, adapted to Palatino mathematics fonts are loaded. The font has a
greater weight and is therefore better suited for high-resolution printing.
Palatino is used in this document and is also set by default in the template.

The font family "`Times"' represents another alternative.
Figure~\ref{fig:tpl:font:newtx} shows a font sample and the corresponding
packages. The sans-serif and thick-weight fonts are identical to Palatino.
